\documentclass[11pt]{article}

% ---------- Packages ----------
\usepackage[a4paper,margin=1in]{geometry}
\usepackage{amsmath,amssymb,amsthm,mathtools}
\usepackage{bbm}
\usepackage{hyperref}
\usepackage{enumitem}
\usepackage{xcolor}

% ---------- Theorem environments ----------
\newtheorem{theorem}{Theorem}
\newtheorem{proposition}{Proposition}
\newtheorem{lemma}{Lemma}
\newtheorem{corollary}{Corollary}

\theoremstyle{definition}
\newtheorem{definition}{Definition}
\newtheorem{remark}{Remark}

% ---------- Commands ----------
\newcommand{\Pcal}{\mathcal{P}}
\newcommand{\Bcal}{\mathcal{B}}
\newcommand{\Fcal}{\mathcal{F}}
\newcommand{\Qcal}{\mathcal{Q}}
\newcommand{\Ccal}{\mathcal{C}}

\pagecolor{white}

% ---------- Title ----------
\title{Observational Foundations of Probability\\
\large A Representation Theorem for Observable Experiments}
\author{Patrick S. Mahon}
\date{}

\begin{document}

\maketitle

\begin{abstract}
Many probabilistic theories begin by postulating an underlying
probability space and then studying the observables derived from it.
This work reverses that perspective.  Starting from observable
experiments organized as a refinement system of queries, we establish
conditions under which compatible observable laws admit a canonical
probability space representation.

We model observations as measurable queries with associated outcome
spaces organized by a refinement preorder. From this system of queries
we construct a canonical realization space as the projective limit of
the observable outcome spaces. The observable $\sigma$-field is then
generated by query-cylinder events.

Within this framework we establish two foundational results. First,
probability measures on the observable $\sigma$-field are uniquely
determined by the distributions of a generating family of queries,
yielding an observational analogue of the classical $\pi$--$\lambda$
uniqueness theorem. Second, compatible families of observable laws
determine a canonical finitely additive observational content on
cylinder events, and this content admits a probabilistic completion
on the realization space by two independent routes.

The \emph{realizability route} (Section~6) takes countably additive
marginals as input and uses a realizability condition on~$\Omega$ to
produce the global measure via a direct Carathéodory argument, with
no topology required on the outcome spaces.  The \emph{Prokhorov route}
(Section~7) starts from only finitely additive marginals and derives
countable additivity from a projective mass-concentration condition
(surjective projective uniform tightness), at the cost of topological
structure on the outcome spaces.  Both routes yield a canonical
probability space $(\Omega,\sigma(\mathcal Q),P)$ whose evaluation
marginals recover the observable laws; under shared hypotheses they
produce the same measure.

The resulting framework connects classical measure-theoretic
probability, statistical experiment theory, and modern observable
approaches to complex systems, and provides a foundation for
extensions involving predictive structure, dynamical observables, and
approximate compatibility.
\end{abstract}

\section{Introduction and motivation}

Probability theory traditionally begins with an underlying probability
space and studies the observables derived from it.  The present work
reverses this perspective by starting from observable experiments
themselves and asking when compatible observable laws admit a canonical
probability space representation.

In many contemporary scientific settings the underlying event algebra
is not directly accessible. Instead, systems are accessed through a
family of admissible measurements, projections, or coarse-grainings.
These observables determine the empirical distinctions that can be
made about the system, while the latent event space remains implicit
or even unknown. This raises a structural question: when can
probabilistic structure be reconstructed from observable experiments
alone?

\paragraph{Problem.}
When do observable experiments determine an underlying probability
space?

\paragraph{Answer.}
We show that compatibility of observable laws is precisely the
condition required for such a representation, yielding a canonical
probability space whose evaluation marginals recover the observable
experiments.  The paper establishes this by two complementary routes.
The realizability route requires: (i) measurable structure on each
query outcome space; (ii) compatible countably additive probability
laws; (iii) directedness conditions; and (iv) realizability of the
projective limit.  The Prokhorov route replaces hypothesis~(ii) by
finitely additive laws satisfying surjective projective uniform
tightness, and derives $\sigma$-additivity from topological
compactness of the outcome spaces.  These two routes resolve the
extension problem from opposite directions: one assumes countable
coherence, the other derives it.

\paragraph{Program perspective.}
This work is part of a broader program developing dynamical structure
from observable experiments.  The present paper establishes the first
stage: the representation theorem.

Classically, Kolmogorov's extension theorem constructs a probability
measure from consistent finite-dimensional marginals.  The present
construction plays an analogous role for observable experiments: the
query system determines a projective family of observable laws whose
compatibility yields the canonical experiment $(\Omega,\mathcal F,P)$.

The structure of the framework is summarized by the following
three-layer architecture.

\begin{center}
\[
\begin{array}{ccc}
\textbf{observational structure} & & \textbf{probabilistic structure} \\[10pt]

\text{queries } (\Qcal,\preceq)
& & \\[6pt]

\downarrow
& & \\[6pt]

\text{projective system of outcomes } \{O_Q,\pi^{Q_2}_{Q_1}\}
& & \text{compatible observable laws } \{\nu_Q\} \\[6pt]

\downarrow
& & \downarrow \\[6pt]

\text{realization space }
\Omega=\varprojlim O_Q
& & \textbf{finite observational content} \\[6pt]

\downarrow
& & \downarrow \\[6pt]

\text{observable }\sigma\text{-field } \sigma(\Qcal)
& & \text{probabilistic completion } P
\end{array}
\]
\end{center}

\paragraph{Main theorem.}
The principal result of this work establishes that compatible families
of observable experiments admit a canonical probabilistic
representation.

\begin{theorem}[Observational representation theorem]
Let $(\Qcal,\preceq)$ be a query system equipped with a compatible
family of observable laws (under the structural hypotheses of
Sections~3 and~6: measurable outcome spaces, countably additive
laws, directedness, and realizability).  Then there exists a
canonical probability space
\[
(\Omega,\mathcal F,P)
\]
such that for every query $Q\in\Qcal$ the evaluation map
\[
Q:\Omega\to O_Q
\]
recovers the corresponding observable law.
\end{theorem}

\paragraph{Relation to the predictive sequel.}
The canonical experiment constructed here provides the probabilistic
representation of observable experiments.  A subsequent work studies
how predictive state spaces and operator dynamics arise naturally from
this observational representation.

\paragraph{Scope.}
The present paper does not derive local measurable structure or local
countable additivity from observation alone. Rather, it shows that
once queries are equipped with measurable outcome spaces and compatible
countably additive observable laws, the remaining global probabilistic
structure is canonically forced under explicit structural assumptions
on the query system. The derivation of local probabilistic structure
from more primitive observable distinctions is a direction for future
work.

\paragraph{Structure of the paper.}
Section~2 recalls the classical uniqueness and extension principles
that motivate the construction.  Section~3 introduces query systems,
their realization spaces, and the observable $\sigma$-field, together
with the three structural conditions used in the main results.
Section~4 proves the observational determination theorem.
Section~5 constructs the finite observational content from compatible
observable laws.  Section~6 proves the observational extension theorem
by a direct Carathéodory argument, establishing the canonical probability
space $(\Omega,\sigma(\Qcal),P)$ without topological assumptions on the
outcome spaces.  Section~7 establishes the Prokhorov extension theorem:
a parallel route in which $\sigma$-additivity is derived rather than
assumed, under surjective projective uniform tightness of the marginals.
The final sections interpret the framework in relation to statistical
experiment theory and outline directions for predictive and dynamical
extensions.

\section{Classical uniqueness and extension}

A central structural feature of classical probability theory is that a probability measure is determined by its values on a generating class of events. One of the most widely used tools expressing this fact is the $\pi$--$\lambda$ theorem~\cite{Billingsley,Kallenberg}. Let $(\Omega,\mathcal F)$ be a measurable space, and let $\mathcal C \subset \mathcal F$ be a $\pi$-system, that is, a collection of sets closed under finite intersections. If $\mathcal C$ generates $\mathcal F$, then any two probability measures $P$ and $P'$ satisfying
\[
P(A) = P'(A), \qquad A \in \mathcal C,
\]
necessarily agree on the full $\sigma$-algebra:
\[
P = P' \text{ on } \mathcal F.
\]
In this sense, the global probabilistic structure is determined by its values on a suitable generating class.

A particularly important application of this principle is Kolmogorov's extension theorem~\cite{Kolmogorov,Billingsley}. In the classical construction of stochastic processes, one begins by specifying a single-time state space $(S,\mathcal S)$ together with an index set $T$, typically representing time. The trajectory space is then taken to be the product
\[
\Omega = S^T,
\]
equipped with the product $\sigma$-algebra
\[
\mathcal F = \mathcal S^{\otimes T}.
\]
The measurable structure is generated by the coordinate projections
\[
\pi_t : \Omega \to S, 
\qquad 
\pi_t(\omega) = \omega(t).
\]
More generally, for finite subsets $\{t_1,\dots,t_n\} \subset T$, one considers the joint projections
\[
\pi_{t_1,\dots,t_n}(\omega)
=
\big(\omega(t_1),\dots,\omega(t_n)\big),
\]
which generate the cylinder $\sigma$-field.

A stochastic process is then specified by a compatible family of finite-dimensional distributions
\[
\mu_{t_1,\dots,t_n}
\in
\mathcal P(S^n),
\]
satisfying the marginal consistency condition
\[
(\pi_{t_1,\dots,t_n})_\# \mu_{t_1,\dots,t_m}
=
\mu_{t_1,\dots,t_n},
\qquad
n < m.
\]
Kolmogorov's theorem asserts that such a compatible family determines a unique probability measure
\[
P \in \mathcal P(\Omega,\mathcal F)
\]
whose finite-dimensional marginals coincide with the prescribed laws.

This construction makes a specific structural assumption about the organization of observational information. The starting point is a global trajectory space together with a canonical system of coordinate projections. Measurable events are generated by restrictions of these coordinates, and probabilistic structure is built relative to this coordinate representation.

However, in many contemporary scientific and statistical settings, this assumption is neither natural nor directly available. Systems are typically accessed through families of admissible measurements, projections, or coarse-grainings, without prior knowledge of a global coordinate structure. The observable interface is therefore given by measurable maps rather than by coordinates of an assumed latent space.

This observation suggests a reformulation in which admissible observations are taken as primitive. Instead of postulating a measurable event space in advance, we derive the observable $\sigma$-field from a family of measurable maps and formulate probabilistic structure in terms of compatibility of their induced laws.

\section{Observational generators}

Replacing coordinates by a more general notion of observable structure
requires specifying the observational interface directly.
Rather than beginning with a fixed measurable state space,
we start from a family of admissible queries.

\begin{definition}[Query]
\label{def:query}
A \emph{query} $Q$ consists of a measurable outcome space
\[
(O_Q,\mathcal B(O_Q)).
\]
\end{definition}

\begin{remark}[Primitive structure]
In the present paper the measurable structure $\mathcal B(O_Q)$ on
each outcome space is taken as a primitive input. It represents the
class of events distinguishable at resolution $Q$. A more foundational
treatment would derive this $\sigma$-algebra from a designated class
$\mathcal E_Q$ of admissible observable distinctions via
$\mathcal B(O_Q) = \sigma(\mathcal E_Q)$; we regard this as a
direction for future work.
\end{remark}

A family of queries $\mathcal Q$ represents the admissible observational
interface through which the system is accessed.

\begin{definition}[Query system]
\label{def:query-system}
A \emph{query system} consists of:
\begin{enumerate}
\item a preorder $(\Qcal,\preceq)$;
\item for each $Q\in\Qcal$, a measurable outcome space
\[
(O_Q,\mathcal B(O_Q));
\]
\item for each relation $Q_1\preceq Q_2$, a measurable refinement map
\[
\pi^{Q_2}_{Q_1}:O_{Q_2}\to O_{Q_1},
\]
satisfying the coherence relations
\[
\pi^Q_Q=\mathrm{id}_{O_Q},
\qquad
\pi^{Q_2}_{Q_1}\circ\pi^{Q_3}_{Q_2}=\pi^{Q_3}_{Q_1}
\quad
(Q_1\preceq Q_2\preceq Q_3).
\]
\end{enumerate}
\end{definition}

In this way a query system $(\Qcal,\preceq)$ may be interpreted as an
information order in the sense of Blackwell and Le Cam~\cite{Blackwell,LeCam}:
queries higher in the order are at least as informative as those below
them, since the outcomes of coarser queries are obtained from finer ones
by measurable post-processing.  Thus the refinement maps are coarsening
maps: whenever $Q_1 \preceq Q_2$, the more informative query $Q_2$
determines the outcome of the less informative query $Q_1$.

We will use two structural hypotheses, depending on the goal.

\begin{definition}[Common coarsenings / lower-directedness]
\label{def:lower-directed}
We say $(\Qcal,\preceq)$ admits \emph{common coarsenings} if for any $Q_1,Q_2\in\Qcal$
there exists $Q_3\in\Qcal$ such that
\[
Q_3 \preceq Q_1
\qquad\text{and}\qquad
Q_3 \preceq Q_2 .
\]
\end{definition}

Intuitively, $Q_3$ is a single (less informative) query whose outcomes can be
pushed forward to recover the outcomes of both $Q_1$ and $Q_2$. This
is the assumption that any two observational resolutions admit a joint
coarser representation.

\begin{definition}[Common refinements / upper-directedness]
\label{def:upper-directed}
We say $(\Qcal,\preceq)$ admits \emph{common refinements} if for any $Q_1,Q_2\in\Qcal$
there exists $Q_3\in\Qcal$ such that
\[
Q_1 \preceq Q_3
\qquad\text{and}\qquad
Q_2 \preceq Q_3 .
\]
\end{definition}

This is the assumption that any two observational resolutions can be
jointly refined by a third --- an assumption about the richness of the
observational interface.

\begin{definition}[Sequential common refinements]
\label{def:seq-upper-directed}
We say $(\Qcal,\preceq)$ admits \emph{sequential common refinements} if every
countable family $Q_1, Q_2, \ldots \in \Qcal$ admits a common refinement: there exists
$Q_*\in\Qcal$ such that
\[
Q_n \preceq Q_*
\qquad\text{for all } n \ge 1.
\]
\end{definition}

This is the assumption that the observational interface is closed
under countable joint refinement: however many query resolutions are
specified, a single finer resolution exists that subsumes all of them.
This is a substantive condition on the interface structure, not merely
a technical regularity hypothesis.

Sequential common refinements imply common refinements (take the pair $Q_1,Q_2$ as
the sequence).  The sequential condition is strictly stronger and is the key analytic
hypothesis enabling $\sigma$-subadditivity of the cylinder premeasure.

\begin{definition}[Realizability]
\label{def:eval-surjective}
The query system is \emph{realizable} if every local outcome extends to a global
realization: for every $Q\in\Qcal$ and every $o\in O_Q$, there exists
$\omega\in\Omega$ such that $\mathrm{eval}_Q(\omega)=o$.

Equivalently, the evaluation maps $\mathrm{eval}_Q:\Omega\to O_Q$ are all surjective.
\end{definition}

\begin{remark}[Status of realizability]
Realizability is a substantive assumption. It asserts that the
projective limit $\Omega$ is non-degenerate in a strong pointwise
sense: every locally consistent outcome extends to a globally coherent
realization. This is used in an essential way in the premeasure
well-definedness argument and in the $\sigma$-subadditivity step of
the extension theorem --- it is not a side condition. Without it,
both steps fail.

A natural measure-theoretic weakening --- likely sufficient for the
proof --- is \emph{almost-sure realizability}: for each $Q$ and
$\nu_Q$-almost every $o \in O_Q$, there exists $\omega \in \Omega$
with $\mathrm{eval}_Q(\omega) = o$. We retain the stronger pointwise
form here for clarity. A sufficient condition for full realizability
is given in Corollary~6.1: when each $O_Q$ is Polish, refinement maps
are continuous, and the index set is countable, the projective limit
is Polish and realizability holds automatically.
\end{remark}

Both structural hypotheses are used in the main results.
Common coarsenings are used for the uniqueness theorem
(the cylinder family forms a $\pi$-system).
Common refinements are used for the extension theorem
(finite cylinder families can be compressed to a single query level
to define a well-posed premeasure).
A sequential strengthening of common refinements is the key analytic
condition that allows the premeasure to be extended to a full
probability measure.

The space of realizations is defined as the set of coherent answers
to all queries:
\[
\Omega
=
\left\{
x \in \prod_{Q\in\mathcal Q} O_Q
\;\middle|\;
x_{Q_1}
=
\pi^{Q_2}_{Q_1}(x_{Q_2})
\text{ whenever } Q_1 \preceq Q_2
\right\}.
\]

Equivalently,
\[
\Omega = \varprojlim_{Q\in\mathcal Q} O_Q ,
\]
the projective limit of the system of outcome spaces with bonding maps
$\pi^{Q_2}_{Q_1}$ \cite{Parthasarathy}.

\begin{remark}[Nondegeneracy of $\Omega$]
The realization space $\Omega$ is a derived object: it is determined
entirely by the query system and carries no structure beyond what the
refinement maps force. It is not assumed as a latent state space.
However, $\Omega$ can be empty or degenerate without additional
hypotheses on the query system --- realizability
(Definition~\ref{def:eval-surjective}) is precisely the condition
asserting non-degeneracy.
\end{remark}

The realization space $\Omega$ serves as a canonical space of coherent
answers to all admissible queries. Probability measures on $\Omega$
represent probabilistic models compatible with the observable interface.

Each query $Q$ induces an evaluation map
\[
\mathrm{eval}_Q : \Omega \to O_Q, \qquad \mathrm{eval}_Q(x) = x_Q.
\]
Probability measures on $\Omega$ represent observational laws compatible
with all admissible queries; the pushforwards $(\mathrm{eval}_Q)_\#P$
recover the observable marginal distributions.

The distribution of a query under a probability measure
$P \in \mathcal P(\Omega,\sigma(\mathcal Q))$
is given by the pushforward measure
\[
Q_\# P
=
(\mathrm{eval}_Q)_\# P
\in
\mathcal P(O_Q).
\]

\begin{lemma}[Projection compatibility]
\label{lem:proj-compat}
For every relation $Q_1 \preceq Q_2$ in $\mathcal Q$ the evaluation maps
satisfy
\[
\mathrm{eval}_{Q_1}
=
\pi^{Q_2}_{Q_1}
\circ
\mathrm{eval}_{Q_2}.
\]
\end{lemma}

\begin{proof}
For $x\in\Omega$ we have
\[
(\pi^{Q_2}_{Q_1}\circ\mathrm{eval}_{Q_2})(x)
=
\pi^{Q_2}_{Q_1}(x_{Q_2}).
\]
By the defining compatibility condition of $\Omega$,
\[
x_{Q_1}=\pi^{Q_2}_{Q_1}(x_{Q_2}),
\]
so the right-hand side equals $\mathrm{eval}_{Q_1}(x)$.
\end{proof}

The observable $\sigma$-field is defined as the smallest
$\sigma$-algebra making all evaluation maps measurable.
Equivalently, it is generated by the cylinder sets
\[
\mathrm{Cyl}(Q,A)
=
\{x\in\Omega : \mathrm{eval}_Q(x)\in A\},
\qquad
A\in\mathcal B(O_Q).
\]

We denote this $\sigma$-field by
\[
\sigma(\mathcal Q).
\]

\begin{lemma}[Measurability of evaluation maps]
\label{lem:eval-measurable}
For every query $Q\in\mathcal Q$, the evaluation map
\[
\mathrm{eval}_Q:\Omega\to O_Q
\]
is measurable with respect to $(\Omega,\sigma(\mathcal Q))$.
\end{lemma}

\begin{proof}
By definition $\sigma(\mathcal Q)$ is the smallest
$\sigma$-algebra making all evaluation maps measurable.
\end{proof}

The common-coarsening hypothesis introduced above is a meet-like
property of the information order on queries. Intuitively, it ensures
that finite collections of observable constraints can be represented at
a single query level.

\begin{lemma}[Cylinder $\pi$-system]
\label{lem:cyl-pisystem}
Assume $(\Qcal,\preceq)$ admits common coarsenings in the sense of
Definition~\ref{def:lower-directed}.

Then the family of cylinder sets
\[
\Ccal_{\Qcal}
=
\{\mathrm{Cyl}(Q,A) : Q\in\mathcal Q,\; A\in\mathcal B(O_Q)\}
\]
forms a $\pi$-system.
\end{lemma}

\begin{proof}
Let $\mathrm{Cyl}(Q_1,A_1)$ and $\mathrm{Cyl}(Q_2,A_2)$ be two cylinders.
The common-coarsening hypothesis provides a query $Q_3$ such that
\[
Q_3 \preceq Q_1
\qquad\text{and}\qquad
Q_3 \preceq Q_2 .
\]
Using the refinement maps
$\pi^{Q_1}_{Q_3}$ and $\pi^{Q_2}_{Q_3}$, a direct calculation shows
\[
\mathrm{Cyl}(Q_1,A_1)\cap\mathrm{Cyl}(Q_2,A_2)
=
\mathrm{Cyl}\!\left(
Q_3,
(\pi^{Q_1}_{Q_3})^{-1}(A_1)
\cap
(\pi^{Q_2}_{Q_3})^{-1}(A_2)
\right).
\]
Since measurable sets are stable under preimages and intersections,
the right-hand side is again a cylinder.
\end{proof}

\begin{remark}[Relation to classical coordinate constructions]
In the classical Kolmogorov construction of stochastic processes, the
indexing family of coordinate projections is ordered by inclusion of
finite coordinate sets and is therefore upper-directed. The
observational framework uses the opposite side of the information order:
common coarsenings permit finite collections of observable constraints
to be represented at a single query level.
\end{remark}


\section{Observational determination}

The following result is the observational analogue of the classical
$\pi$--$\lambda$ uniqueness theorem.

\begin{theorem}
\label{thm:obs-determination}
Let $(\Qcal,\preceq)$ admit common coarsenings (Definition~\ref{def:lower-directed}) and let
$P$ and $P'$ be probability measures on the observable space
$(\Omega,\sigma(\Qcal))$ generated by $\mathcal Q$.

If
\[
Q_\# P = Q_\# P'
\]
for all $Q \in \Qcal$, then $P = P'$.
\end{theorem}

\begin{proof}
Equality of pushforward laws implies that $P$ and $P'$ agree on all
cylinder events.
By Lemma~\ref{lem:cyl-pisystem} the cylinder family forms a $\pi$-system generating
$\sigma(\Qcal)$.
The classical $\pi$--$\lambda$ theorem therefore implies that the two
measures coincide on the entire observable $\sigma$-field.
\end{proof}

\begin{remark}
The determination theorem is purely measure-theoretic.  It depends
only on the observable $\sigma$-field generated by queries and on
the lower-directedness hypothesis ensuring that cylinder events form
a $\pi$-system.  No compactness or topological assumptions enter
at this stage.
\end{remark}

\section{Finite observational content}

\begin{definition}[Cylinder set semiring]
Let
\[
\Ccal_{\Qcal}
=
\{\mathrm{Cyl}(Q,A): Q\in\Qcal,\;A\in\mathcal B(O_Q)\}
\]
denote the family of cylinder sets.
\end{definition}

\begin{lemma}[Cylinder set semiring structure]
\label{lem:cyl-semiring}
Assume $(\Qcal,\preceq)$ admits common coarsenings
(Definition~\ref{def:lower-directed}).
Then $\Ccal_{\Qcal}$ is a set semiring: it is closed under finite
intersections (already the $\pi$-system property of
Lemma~\ref{lem:cyl-pisystem}) and set differences decompose as finite
disjoint unions of cylinders.
\end{lemma}

\begin{proof}[Proof sketch]
Intersection: Lemma~\ref{lem:cyl-pisystem}.
Difference: given $\mathrm{Cyl}(Q_1,A)\setminus\mathrm{Cyl}(Q_2,B)$,
use a common coarsening $Q_3 \preceq Q_1, Q_2$ to write both as
cylinders at $Q_3$; the difference is then a single cylinder at $Q_3$
and is itself in $\Ccal_{\Qcal}$.
\end{proof}

\begin{lemma}[Finite cylinder compression]
\label{lem:finite-cylinder}
Assume $(\Qcal,\preceq)$ admits common refinements
(Definition~\ref{def:upper-directed}).

Then for any finite family of queries $Q_1,\ldots,Q_n$ and measurable
sets $A_i\in\mathcal B(O_{Q_i})$, there exists a common refinement
$Q_*$ such that
\[
\bigcap_{i=1}^n \mathrm{Cyl}(Q_i,A_i)
=
\mathrm{Cyl}\!\left(Q_*,\;\bigcap_{i=1}^n (\pi^{Q_*}_{Q_i})^{-1}(A_i)\right).
\]
\end{lemma}

\begin{proof}
Upper-directedness gives a common refinement $Q_*$ with $Q_i\preceq Q_*$
for all $i$.  Then $\pi^{Q_*}_{Q_i}:O_{Q_*}\to O_{Q_i}$ is well-typed
and $\mathrm{Cyl}(Q_i,A_i)=\mathrm{Cyl}(Q_*,(\pi^{Q_*}_{Q_i})^{-1}(A_i))$
by the refinement identity.  The intersection follows.
\end{proof}

\begin{remark}
Lemma~\ref{lem:finite-cylinder} uses common \emph{refinements} rather
than common coarsenings.  The direction matters: the projection
$\pi^{Q_*}_{Q_i}:O_{Q_*}\to O_{Q_i}$ goes from the finer space to the
coarser one, so it is the preimage that lives in $O_{Q_*}$ and the
constraint set is well-typed.  An upper bound in the information order
is what is needed here.
\end{remark}

We now turn to the question of probabilistic realization.
Suppose we are given, for each query $Q \in \mathcal Q$, a probability law
\[
\nu_Q \in \mathcal P(O_Q).
\]

\begin{definition}[Observational compatibility]
\label{def:obs-compat}
A family $\{\nu_Q\}_{Q \in \mathcal Q}$ with $\nu_Q \in \mathcal P(O_Q)$
is \emph{compatible} if for every relation $Q_1 \preceq Q_2$ in the query preorder,
\[
(\pi^{Q_2}_{Q_1})_\# \nu_{Q_2}=\nu_{Q_1}.
\]
\end{definition}

In other words, whenever one query is obtained from another by
coarsening, the corresponding probability laws agree under the induced
marginalization.

\begin{theorem}[Finite observational content]
\label{thm:finite-content}
Suppose:
\begin{enumerate}
\item The refinement preorder on $\Qcal$ admits common coarsenings
(Definition~\ref{def:lower-directed}) and common refinements
(Definition~\ref{def:upper-directed}).
\item The query system is realizable (Definition~\ref{def:eval-surjective}).
\item The family $\{\nu_Q\}_{Q \in \mathcal Q}$ is observationally compatible
in the sense of Definition~\ref{def:obs-compat}.
\end{enumerate}
Then the family $\{\nu_Q\}$ determines a well-defined finitely
additive content $\mu_0$ on $\Ccal_{\Qcal}$:
\[
\mu_0(\mathrm{Cyl}(Q,A)) := \nu_Q(A), \qquad A\in\mathcal B(O_Q).
\]
\end{theorem}

\begin{proof}[Proof sketch]
\textit{Well-definedness.}
Suppose $\mathrm{Cyl}(Q,A)=\mathrm{Cyl}(Q',A')$.  Use common refinements
to find $Q_*$ finer than both $Q$ and $Q'$.  The constraint sets
at $Q_*$ are the preimages of $A$ and $A'$ under $\pi^{Q_*}_Q$ and
$\pi^{Q_*}_{Q'}$ respectively.  Since $\mathrm{Cyl}(Q,A)=\mathrm{Cyl}(Q',A')$,
surjectivity of $\mathrm{eval}_{Q_*}$ (realizability) implies the two
preimage sets in $O_{Q_*}$ coincide.  Compatibility then gives
$\nu_Q(A)=\nu_{Q_*}(\text{preimage})=\nu_{Q'}(A')$.

\textit{Finite additivity.}
For disjoint cylinders with measurable union, the same compression to
a common refinement level $Q_*$ reduces additivity to that of $\nu_{Q_*}$.
\end{proof}

\begin{remark}
The passage from finite content to a countably additive
probability measure on $(\Omega,\sigma(\Qcal))$ requires one further
condition: $\sigma$-subadditivity of $\mu_0$.  This is the analytic
heart of the extension problem, and the next section shows how
sequential common refinements supply it.

Hypothesis~(3) requires each $\nu_Q$ to be a countably additive
probability measure. This is assumed, not derived: the countable
additivity present in the global extension theorem (\S6) is propagated
from this local assumption via the compression argument, not
constructed from below. The question of whether countable additivity
of the $\nu_Q$ can itself be derived from compatibility and interface
structure alone is a direction for future work.
\end{remark}

\section{Observational extension theorem}

The finite content $\mu_0$ constructed in the previous section is
finitely additive on the cylinder set semiring $\Ccal_{\Qcal}$.
The central result of this paper shows that sequential common refinements
supply the missing $\sigma$-subadditivity, allowing Carathéodory's
extension theorem~\cite{Billingsley,Bogachev} to produce a full probability
measure on $(\Omega,\sigma(\Qcal))$.

\begin{theorem}[Observational extension theorem]
\label{thm:obs-extension}
Suppose:
\begin{enumerate}
\item $(\Qcal,\preceq)$ admits common coarsenings
      (Definition~\ref{def:lower-directed}).
\item $(\Qcal,\preceq)$ admits sequential common refinements
      (Definition~\ref{def:seq-upper-directed}).
\item The query system is realizable
      (Definition~\ref{def:eval-surjective}).
\item The family $\{\nu_Q\}_{Q\in\Qcal}$ is observationally compatible
      (Definition~\ref{def:obs-compat}) and each $\nu_Q$ is a
      probability measure.
\end{enumerate}
Then there exists a unique probability measure
\[
P \in \mathcal P(\Omega,\sigma(\Qcal))
\]
such that for every $Q\in\Qcal$,
\[
(\mathrm{eval}_Q)_\# P = \nu_Q.
\]
\end{theorem}

\begin{proof}
\textit{Step 1: Finite content.}
Theorem~\ref{thm:finite-content} gives a well-defined finitely additive
content $\mu_0$ on $\Ccal_{\Qcal}$ with $\mu_0(\mathrm{Cyl}(Q,A))=\nu_Q(A)$.
By Lemma~\ref{lem:cyl-semiring}, $\Ccal_{\Qcal}$ is a set semiring.

\medskip
\noindent\textit{Step 2: $\sigma$-subadditivity.}
Let $\mathrm{Cyl}(Q,B)\subseteq\bigcup_{n\ge1}\mathrm{Cyl}(Q_n,A_n)$ with all
sets in $\Ccal_{\Qcal}$.  Apply sequential common refinements to the
sequence $Q,Q_1,Q_2,\ldots$ to obtain a single query $Q_*$ with
$Q\preceq Q_*$ and $Q_n\preceq Q_*$ for all $n$.

Compression (Lemma~\ref{lem:finite-cylinder}) rewrites each cylinder
at level $Q_*$:
\[
\mathrm{Cyl}(Q,B) = \mathrm{Cyl}(Q_*,E),
\qquad
\mathrm{Cyl}(Q_n,A_n) = \mathrm{Cyl}(Q_*,E_n),
\]
where $E=(\pi^{Q_*}_Q)^{-1}(B)$ and $E_n=(\pi^{Q_*}_{Q_n})^{-1}(A_n)$
are measurable subsets of $O_{Q_*}$.
Realizability and the set inclusion in $\Omega$ imply $E\subseteq\bigcup_n E_n$ in $O_{Q_*}$.
Therefore
\[
\mu_0(\mathrm{Cyl}(Q,B))
= \nu_{Q_*}(E)
\le \sum_{n=1}^\infty \nu_{Q_*}(E_n)
= \sum_{n=1}^\infty \mu_0(\mathrm{Cyl}(Q_n,A_n)),
\]
where the inequality is $\sigma$-subadditivity of the single measure $\nu_{Q_*}$.

\medskip
\noindent\textit{Step 3: Carathéodory extension.}
Since $\Ccal_{\Qcal}$ is a set semiring generating $\sigma(\Qcal)$ and
$\mu_0$ is a $\sigma$-subadditive additive content on it,
Carathéodory's extension theorem~\cite{Billingsley,Bogachev} applies
and yields a measure $P$ on $(\Omega,\sigma(\Qcal))$ extending $\mu_0$.

\medskip
\noindent\textit{Step 4: Marginal recovery.}
For any $A\in\mathcal B(O_Q)$,
\[
P(\mathrm{eval}_Q^{-1}(A))
= P(\mathrm{Cyl}(Q,A))
= \mu_0(\mathrm{Cyl}(Q,A))
= \nu_Q(A),
\]
so $(\mathrm{eval}_Q)_\# P = \nu_Q$.

\medskip
\noindent\textit{Step 5: $P$ is a probability measure.}
$P(\Omega)=P(\mathrm{Cyl}(Q,O_Q))=\nu_Q(O_Q)=1$ for any $Q$.

\medskip
\noindent\textit{Uniqueness.}
Two probability measures with the same evaluation marginals agree on
every cylinder event, hence on $\sigma(\Qcal)$ by
Theorem~\ref{thm:obs-determination}.
\end{proof}

\begin{remark}[Assumptions versus derived structure]
The theorem assumes: (i) measurable structure on each outcome space;
(ii) countable additivity of each observable law $\nu_Q$; (iii)
lower-directedness, sequential upper-directedness, and realizability
of the query system. It derives: a canonical $\sigma$-additive
probability measure on $(\Omega,\sigma(\Qcal))$ whose evaluation
marginals recover all observable laws.

Countable additivity of $P$ is not constructed from first principles
but propagated from the countable additivity of $\nu_{Q_*}$ via the
compression argument. Assumptions (i) and (ii) are inputs from
classical measure theory; a more foundational treatment would aim to
derive them from the structure of observable distinctions. Assumptions
(iii) are substantive conditions on the query system: they express
joint coarsenability, closure under countable joint refinement, and
non-degeneracy of the realization space.
\end{remark}

\begin{remark}[Structure of the proof]
The proof uses the two directedness conditions for distinct roles.
Common coarsenings give the semiring structure on $\Ccal_{\Qcal}$
(needed for Carathéodory) and the $\pi$-system structure needed for
uniqueness.  Sequential common refinements compress any countable
cylinder family to a single query level, reducing $\sigma$-subadditivity
of $\mu_0$ to $\sigma$-subadditivity of the single measure $\nu_{Q_*}$,
which is immediate.  Realizability is used only to translate set
inclusions in $\Omega$ to set inclusions in $O_{Q_*}$.
\end{remark}

\begin{remark}[Relation to Kolmogorov's extension theorem]
Kolmogorov's theorem constructs a probability measure on $S^T$ from
consistent finite-dimensional distributions.  The index set is a
coordinate product and the consistency condition is the push-forward
condition under coordinate projections.

The observational extension theorem plays the same role for query
systems.  The realization space $\Omega=\varprojlim O_Q$ replaces the
product $S^T$, compatible marginals replace consistent
finite-dimensional distributions, and sequential common refinements
replace the countable-cofinality condition that reduces the Kolmogorov
argument to a countable system.

The key conceptual difference is that the observational route imposes
no topology on the outcome spaces.  The entire construction is
measure-theoretic: the measure $P$ is built directly on the
$\sigma(\Qcal)$-measurable structure of $\Omega$ by Carathéodory's
theorem, without any appeal to tightness or compactness.
\end{remark}

\begin{corollary}[Topological setting]
\label{cor:topological}
In particular, Corollary~\ref{cor:topological} provides the main
route to verifying realizability: under Polish outcome spaces and a
countable cofinal chain, all four hypotheses of
Theorem~\ref{thm:obs-extension} are automatically satisfied.

The hypotheses of Theorem~\ref{thm:obs-extension} are satisfied
whenever:
\begin{enumerate}
\item each outcome space $O_Q$ is Polish and the refinement maps are
      continuous (so $\Omega$ is a closed subspace of a countable
      product of Polish spaces, hence Polish~\cite{Bogachev} and realizable);
\item the query preorder contains a countable cofinal chain
      $Q_1\preceq Q_2\preceq\cdots$ (so every sequence of queries
      admits a common refinement within this chain).
\end{enumerate}
In this case the compatible observable laws need not be assumed Radon:
the measure $P$ produced by Theorem~\ref{thm:obs-extension} is
automatically a Radon probability measure on the Polish space $\Omega$.
\end{corollary}

\begin{remark}[Lean formalization]
\label{rem:lean-caratheodory}
The proof of Theorem~\ref{thm:obs-extension} given above is the proof
that is fully formalized in Lean~4~\cite{Lean4} (\texttt{observational\_extension}
in \texttt{QuerySystem.lean}, 0 sorrys).
The Lean hypotheses correspond exactly to the four conditions of the
theorem: \texttt{LowerDirected} (common coarsenings),
\texttt{SequentiallyUpperDirected} (sequential common refinements),
\texttt{EvalSurjective} (realizability), and
\texttt{CompatibleMarginals} with \texttt{IsProbabilityMeasure}.
The Carathéodory step uses \texttt{AddContent.measure} from Mathlib~\cite{Mathlib}.
\end{remark}

\section{Prokhorov extension for query systems}
\label{sec:prokhorov}

Theorem~\ref{thm:obs-extension} proceeds from the \emph{realizability} side:
given $\sigma$-additive marginals $\{\nu_Q\}$ and a realizability
condition on~$\Omega$, it produces the global measure~$P$.
This section establishes a route from the opposite direction, the \emph{tightness} side:
given only \emph{finitely} additive marginals, a projective mass-concentration
condition forces $\sigma$-additivity and delivers $P$ directly.
These two theorems resolve the same extension problem from opposite directions:
one assumes countable coherence, the other derives it.

\subsection*{Projective uniform tightness}

Because the marginals $\{\nu_Q\}$ live on different spaces $O_Q$,
the classical Prokhorov notion of tightness cannot be applied
level by level.  The correct formulation is projective: compact
approximants must be chosen compatibly across levels so that they
assemble into a compact subset of the inverse limit~$\Omega$.

\begin{definition}[Surjective projective uniform tightness]
\label{def:proj-tight}
A finitely compatible family $\{\nu_Q\}_{Q\in\Qcal}$ of finitely
additive normalized measures on $(\mathcal{B}(O_Q))_{Q\in\Qcal}$
is \emph{surjectively projectively uniformly tight} if for every
$\varepsilon>0$ there exists a system $\{K_Q\}_{Q\in\Qcal}$ satisfying:
\begin{enumerate}
\item $K_Q\subset O_Q$ is compact for each $Q\in\Qcal$,
\item $\pi^{Q_2}_{Q_1}(K_{Q_2})=K_{Q_1}$ for all $Q_1\preceq Q_2$
      \emph{(surjectivity of bonding maps on compact sets)},
\item $\nu_Q(O_Q\setminus K_Q)<\varepsilon$ for every $Q\in\Qcal$.
\end{enumerate}
Condition~(2) ensures the compact sets form a consistent inverse system,
so their intersection in $\Omega$,
\[
K_\Omega := \Omega\cap\prod_{Q\in\Qcal} K_Q = \varprojlim K_Q,
\]
is nonempty and compact, and each projection
$\mathrm{eval}_Q:K_\Omega\to K_Q$ is surjective (by the
inverse-limit theorem for compact Hausdorff spaces with surjective
bonding maps).
\end{definition}

\begin{remark}[Relation to individual tightness]
Definition~\ref{def:proj-tight} strengthens individual tightness
(compact $K_Q$ at each level independently) by requiring the compact
approximants to be chosen compatibly under refinement with surjective
bonding maps.  This cross-level compatibility is precisely what is
needed to lift a local nonemptiness argument into $\Omega$.
\end{remark}

\subsection*{The theorem}

\begin{theorem}[Prokhorov extension for query systems]
\label{thm:prokhorov}
Let $(\Qcal,\preceq)$ be a countable sequentially upper-directed
query system in which each outcome space $O_Q$ is Hausdorff with
$\mathcal{E}_Q=\mathcal{B}(O_Q)$, and each refinement map
$\pi^{Q_2}_{Q_1}:O_{Q_2}\to O_{Q_1}$ is continuous.
Let $\{\nu_Q\}_{Q\in\Qcal}$ be a finitely compatible family of
finitely additive normalized measures on $\mathcal{B}(O_Q)$.

If $\{\nu_Q\}$ is surjectively projectively uniformly tight
(Definition~\ref{def:proj-tight}), then there exists a unique
$\sigma$-additive probability measure $P$ on $(\Omega,\sigma(\Qcal))$
satisfying
\[
(\mathrm{eval}_Q)_\#P=\nu_Q \qquad\text{for every }Q\in\Qcal.
\]
Uniqueness is among all $\sigma$-additive probability measures on
$(\Omega,\sigma(\Qcal))$ with the prescribed marginals.
\end{theorem}

\begin{proof}[Proof sketch]
\textit{Step~1: Compact core.}
For each $\varepsilon>0$ apply Definition~\ref{def:proj-tight} to
obtain $\{K_Q\}$ with surjective bonding maps, compact $K_\Omega=\varprojlim K_Q$,
and $\nu_Q(O_Q\setminus K_Q)<\varepsilon$ for all~$Q$.

\medskip
\noindent\textit{Step~2: Closed-cylinder continuity from above.}
Let $C_N=\mathrm{Cyl}(Q_N,F_N)\downarrow\emptyset$ with each $F_N\subseteq O_{Q_N}$
closed.  Assuming $\mu_0(C_N)\ge\varepsilon$ for all $N$ (seeking contradiction),
the mass estimate $\nu_{Q_N}(F_N\cap K_{Q_N})\ge\varepsilon/2>0$ combined with
surjectivity of $\mathrm{eval}_{Q_N}:K_\Omega\to K_{Q_N}$ gives
$C_N\cap K_\Omega\ne\emptyset$ for all~$N$.
Since $F_N$ is closed and $\mathrm{eval}_{Q_N}$ is continuous, each
$C_N\cap K_\Omega$ is closed in the compact space $K_\Omega$.
The finite intersection property yields
$\bigcap_N(C_N\cap K_\Omega)\ne\emptyset$,
contradicting $C_N\downarrow\emptyset$ in~$\Omega$.

\medskip
\noindent\textit{Step~3: Extension to Borel cylinders.}
Each $K_Q$ is compact Hausdorff, so $\nu_Q|_{K_Q}$ is regular by
Alexandrov's theorem (no additional hypothesis required).
Inner regularity approximates every Borel cylinder by closed cylinders
to within $O(\varepsilon)$.
Sequential upper-directedness provides common refinements for finite
intersections of closed cylinders, preserving closedness.
A standard approximation argument reduces the general Borel case to
the closed-cylinder case already handled.
Carathéodory's theorem then extends the $\sigma$-additive premeasure
$\mu_0$ uniquely to $\sigma(\Ccal_{\Qcal})=\sigma(\Qcal)$ on the
ambient product.

\medskip
\noindent\textit{Step~4: Concentration on $\Omega$ (graph-support lemma).}
For each comparable pair $Q_1\preceq Q_2$, the intersection
$\mathrm{Cyl}(Q_1,A)\cap\mathrm{Cyl}(Q_2,B)=\mathrm{Cyl}(Q_2,(\pi^{Q_2}_{Q_1})^{-1}(A)\cap B)$
shows the joint law of $(\mathrm{eval}_{Q_1},\mathrm{eval}_{Q_2})$ under $P$
equals $(\mathrm{graph}_{\pi})_\#\nu_{Q_2}$, which is supported on
the graph $G_{Q_1,Q_2}=\{(a,b):a=\pi^{Q_2}_{Q_1}(b)\}$.
Hence the defect set $D_{Q_1,Q_2}:=\{x:x_{Q_1}\ne\pi^{Q_2}_{Q_1}(x_{Q_2})\}$
has $P(D_{Q_1,Q_2})=0$.
Since $\Qcal$ is countable there are countably many comparable pairs;
a countable union of null sets gives $P(\Omega)=1$.
Restricting $P$ to $\Omega$ yields the required measure.
Uniqueness follows because cylinder sets generate $\sigma(\Qcal)$
and determine any $\sigma$-additive probability measure with the
given marginals.
\end{proof}

\begin{proposition}[Consistency between routes]
\label{prop:route-consistency}
If $\{\nu_Q\}$ is a family of $\sigma$-additive probability measures
satisfying both the hypotheses of Theorem~\ref{thm:obs-extension}
and those of Theorem~\ref{thm:prokhorov}, then both constructions
produce the same measure~$P$.
\end{proposition}
\begin{proof}
Both measures are $\sigma$-additive on $(\Omega,\sigma(\Qcal))$
with marginals $\nu_Q$ for every~$Q$.
Since cylinder sets generate $\sigma(\Qcal)$ and form a
$\pi$-system, uniqueness of $\sigma$-additive extensions implies
equality.
\end{proof}

\begin{remark}[Comparison of the two routes]
The two extension theorems address complementary scenarios.

\smallskip
\begin{center}
\begin{tabular}{lll}
\hline
Route & Input & Key mechanism \\
\hline
Realizability (Thm.~\ref{thm:obs-extension}) &
  $\sigma$-additive $\nu_Q$ & realizability of $\Omega$ \\
Prokhorov (Thm.~\ref{thm:prokhorov}) &
  finitely additive $\nu_Q$ & projective tightness \\
\hline
\end{tabular}
\end{center}

\smallskip\noindent
The realizability route assumes countable coherence of the marginals and
requires no topology on the outcome spaces.  The Prokhorov route derives
countable coherence from a topological mass-concentration condition:
surjective projective compact cores prevent mass from escaping the
inverse limit.  Both routes conclude that a unique global measure $P$
on $(\Omega,\sigma(\Qcal))$ exists with the prescribed marginals.

In practice, the realizability route applies when countable additivity
of the marginals is known a priori; the Prokhorov route applies when
only finite compatibility is available but surjective projective
uniform tightness (Definition~\ref{def:proj-tight}) can be verified.

The role of projective tightness in the absence of countability, and
its relation to structural conditions on the query system, is discussed
in Section~\ref{sec:outlook}.
\end{remark}

\section{Representation and experiment-theoretic interpretation}

The preceding results show that the probabilistic content of a model is
governed by observable regularities and their compatibility relations.
Compatible observable marginals determine a canonical content on
cylinder events; sequential common refinements allow this content to
be extended to a full probability measure on the observational horizon.
Different realizations may therefore support the same observational
laws, which motivates the following notion.

The structural relation to the theory of statistical experiments may
also be made explicit.  In the sense of Blackwell and Le Cam~\cite{Blackwell,LeCam},
each query together with its induced law defines an observable experiment,
and the refinement preorder records the corresponding informativeness
order under measurable post-processing.  From this perspective, the
present framework and classical experiment theory are built from
closely related ingredients:
\[
\begin{array}{cc}
\textbf{Le Cam experiment theory} & \textbf{observational framework} \\[6pt]
\text{statistical experiment} & \text{query} \\
\text{parameter space} & \text{realization space } \Omega \\
\text{observation map} & \text{evaluation map} \\
\text{post-processing / garbling} & \text{query refinement} \\
\text{experiment family} & \text{query system} \\
\text{experiment comparison} & \text{informativeness preorder} \\
\text{experiment equivalence} & \text{observational equivalence}
\end{array}
\]

The difference lies in the direction of construction. Classical
experiment theory begins with an underlying probabilistic model and
asks how different observational mechanisms compare. The present
framework begins with a compatible family of observable experiments
and asks when they determine a probabilistic model. In this sense,
Le Cam theory moves from models to experiments, whereas the present
framework moves from experiments to models.

More concretely, in classical experiment theory a statistical
experiment is given by a family
\[
\theta \mapsto P_\theta .
\]
In the present setting an observable experiment is given by a query
\[
\omega \mapsto Q(\omega).
\]
The refinement relation
\[
Q_1 \preceq Q_2
\quad\Longleftrightarrow\quad
Q_1 = \pi \circ Q_2
\]
is therefore the same structural relation that appears in the
Blackwell--Le Cam theory of experiment comparison \cite{Blackwell,LeCam,Torgersen}: one experiment is
obtained from another by measurable post-processing.  The query system
may thus be interpreted as a preorder of observable experiments
ordered by informativeness.

\begin{definition}[Observational equivalence]
Two realizations are observationally equivalent if they induce the same
pushforward laws for every admissible query.
\end{definition}

Different measurable realizations may support the same observable laws.
In particular, $\Omega$ is not a hidden state space postulated
independently of the observable interface: it is the canonical
consistency object determined by the query system, carrying exactly
the structure forced by the refinement maps and no more.
In that sense, the latent measurable space functions as a representation
of the observable structure rather than as a primitive object.
From this viewpoint the entire framework may be read as an inverse
problem for experiment theory: instead of comparing experiments
relative to a fixed model, one reconstructs the model from a
compatible family of experiments.

This viewpoint is especially relevant for subsequent developments.
There the basic progression is not merely
\[
\text{queries}
\rightarrow
\text{predictive experiments}
\rightarrow
\text{predictive sufficiency}
\rightarrow
\text{observable operators},
\]
but more specifically a passage from observable experiment structure
to predictive experiment structure.  The broader program may therefore
be summarized as
\[
\text{experiment lattice} \rightarrow \text{observable compatibility} \rightarrow \text{probability},
\]
so that probability appears as a representation theorem for observable
information structures.

The predictive structure that emerges from the canonical experiment
$(\Omega,\sigma(\Qcal),P)$ --- predictive kernels, minimal predictive
state spaces, and operator representations --- is studied in subsequent
work.

\section{Approximate compatibility}

In practical settings, exact compatibility is rarely satisfied. Instead, observational laws are typically consistent only up to controlled discrepancies. A key question is whether control of divergences
\[
D(Q_\# P \,\|\, Q_\# P')
\]
on generators implies global control. This suggests several investigative directions into robustness, optimal transport, and information geometry:
\begin{itemize}
\item Stability and approximate extension under bounds on divergences such as total variation, $f$-divergences, or Wasserstein metrics \cite{Villani}.
\item Information-geometric formulations of observational distinguishability and compatibility.
\item Robust probabilistic modeling in high-dimensional and partially observed systems.
\end{itemize}

These problems are closely related to contemporary themes in statistics and machine learning, including distributional robustness, uncertainty quantification, and inference under model misspecification. The observational perspective suggests that probabilistic structure may be viewed as the most stable completion of empirical regularities rather than as an exact underlying law.

\section{Outlook and connections}
\label{sec:outlook}

The observational formulation developed here provides a unified viewpoint linking several classical and modern strands of probability and statistics:

\begin{enumerate}
	\item Classical stochastic processes arise as a special case in which the observational interface consists of coordinate projections of a product space $S^T$.  The observational extension theorem recovers Kolmogorov’s construction in this setting (via Corollary~\ref{cor:topological}) while clarifying its structural content: the key hypothesis is a countable cofinality condition on the index set, not topology on the factor spaces.

	\item Symmetry and exchangeability constraints on observational laws lead naturally to representation theorems such as de Finetti’s~\cite{deFinetti}. In this setting, mixture representations appear as structural consequences of invariance of the observational interface.

	\item The theory of statistical experiments and sufficiency can be interpreted in terms of refinement and equivalence of observational systems. Minimal sufficient representations correspond to minimal observational structures preserving predictive content.

	\item The framework connects naturally to dynamical and operator-theoretic approaches. In time-evolving systems, observables generate predictive structure, and Markov or semigroup properties may be interpreted as emergent regularities of compatible observational laws. This perspective suggests new connections to transfer operators, Koopman theory \cite{Koopman1931,Mezic}, and data-driven dynamical modeling.

	\item The observational viewpoint provides a potential synthesis of classical objectivist and Bayesian interpretations. Stable empirical regularities and coherent inference both arise as properties of compatible observational systems, without requiring a prior commitment to either ontological or subjective foundations.
\end{enumerate}

These connections motivate further development of approximate, dynamical,
and geometric extensions of the theory. In particular, the emergence of
predictive structure, robustness under partial observability, and
operator-based representations of observational consistency form key
directions for future work.

\paragraph{Progress on structural forcing of $\sigma$-additivity.}
Section~7 derives $\sigma$-additivity from projective tightness in
the countable setting, taking $\sigma$-additive marginals as input.
A companion paper, \emph{Finitely Additive Observable Laws and the
Prokhorov Extension for Query Systems}, shows that the $\sigma$-additivity
assumption on individual marginals can itself be dropped in the compact
setting: when each outcome space $O_Q$ is compact Hausdorff and each
refinement map is surjective and continuous, any finitely compatible
family of normalized valuations extends to a unique $\sigma$-additive
global probability measure.  Countable additivity at the marginal level
is a theorem of system architecture, not an assumption.

The remaining open question concerns the infinite-range case.
For query systems whose outcome spaces are not compact, one requires
a \emph{surjective projective uniform tightness} (SPUT) condition on
the compatible family --- a projective analogue of Prokhorov tightness
with surjectivity of bonding maps on compact cores playing the role
that single-space compactness plays classically.  Whether SPUT can
itself be derived from purely combinatorial or order-theoretic
conditions on $(\Qcal, \preceq)$, without topological hypotheses on
the outcome spaces, remains open.


\bibliographystyle{plain}
\bibliography{references}

\section*{Acknowledgements}

The author thanks Claude (Anthropic) for assistance with mathematical
writing, Lean~4 formalization, and technical discussion throughout the
development of this work.  The author gratefully acknowledges financial
support from Dr.\ Simon Bonner and Dr.\ Matt Davison at the University
of Western Ontario, and from MITACS.

\end{document}